% Generated from validated synthesis. Do not hand-edit.
\section*{Monthly Signals}
\addcontentsline{toc}{section}{Monthly Signals}
\smallskip
\noindent\textbf{Frontier Models} 単純な性能順位より、提供形態、価格、API、open model、agent/codingへの接続という複数の競争軸が前景化した。 \hfill{\footnotesize p.~\pageref{pkg:frontier-models-july}}\par
\smallskip
\noindent\textbf{Multimodal} 生成modalitiesの追加だけでなく、リアルタイムinteraction、cross-modal reference、region別availabilityへ競争が広がった。 \hfill{\footnotesize p.~\pageref{pkg:multimodal-july}}\par
\smallskip
\noindent\textbf{Inference \& Serving} model architectureとhardwareの変化を受け、runner、speculative decoding、MoE kernel、schedulerまでserving stack側の更新が進んだ。 \hfill{\footnotesize p.~\pageref{pkg:inference-serving-july}}\par
\smallskip
\noindent\textbf{長時間エージェントと安全境界} 長時間taskではevaluationとruntimeの双方が変わり、trajectory、external interface、persistent memoryまで安全性の評価対象が広がった。 \hfill{\footnotesize p.~\pageref{pkg:agents-july} / p.~\pageref{pkg:agent-safety-security-july}}\par
\smallskip
\noindent\textbf{Paper Watch} 自己回帰以外のgeneration processを試すopen-weight研究が現れ、servingの前提そのものを変える可能性が観測対象に加わった。 \hfill{\footnotesize p.~\pageref{pkg:paper-watch-july}}\par
\medskip
\begin{claimboundary}[Retrospective scope]
本号は2026年7月を後日確認可能になった一次情報も用いて再構成するRetrospective Specialである。Coverage Periodと制作時点を同一視せず、vendor / project / author claimの境界は各記事とTechnical Notesで明示する。
\end{claimboundary}
\medskip
\tableofcontents
